% !TeX root=../main.tex
% در این فایل، عنوان پایان‌نامه، مشخصات خود، متن تقدیمی‌، ستایش، سپاس‌گزاری و چکیده پایان‌نامه را به فارسی، وارد کنید.
% توجه داشته باشید که جدول حاوی مشخصات پروژه/پایان‌نامه/رساله و همچنین، مشخصات داخل آن، به طور خودکار، درج می‌شود.
%%%%%%%%%%%%%%%%%%%%%%%%%%%%%%%%%%%%
% دانشگاه خود را وارد کنید
\university{دانشگاه تهران}
% پردیس دانشگاهی خود را اگر نیاز است وارد کنید (مثال: فنی، علوم پایه، علوم انسانی و ...)
\college{پردیس دانشکده‌های فنی}
% دانشکده، آموزشکده و یا پژوهشکده  خود را وارد کنید
% TODO تأیید کنید — از قالب رسمی BScThesis_template_new.docx همین پروژه استخراج شده
\faculty{دانشکده مهندسی برق و کامپیوتر}
% گروه آموزشی خود را وارد کنید (در صورت نیاز)
% TODO — نام دقیق گروه را وارد کنید (این مقدار placeholder است تا سند کامپایل شود)
\department{TODO}
% رشته تحصیلی خود را وارد کنید
% TODO تأیید کنید
\subject{مهندسی کامپیوتر}
% گرایش خود را وارد کنید
% TODO — نام دقیق گرایش را وارد کنید (مثلاً هوش مصنوعی؛ placeholder تا تکمیل)
\field{TODO}
% عنوان پایان‌نامه را وارد کنید
% TODO تأیید یا اصلاح کنید — پیشنهاد بر اساس AGENTS.md (تعریف پروژه)
\title{پیش‌بینی نرخ عدم‌دریافت غذا در سلف‌های دانشگاه تهران با رگرسیون کوانتایل روی
سری‌های زمانی چندگانه}
% نام استاد(ان) راهنما را وارد کنید
% TODO — نام و مرتبه‌ی استاد راهنما را وارد کنید
\firstsupervisor{}
\firstsupervisorrank{}
% اگر استاد راهنمای دوم ندارید، دو خط زیر را غیرفعال کنید.
%\secondsupervisor{دکتر راهنمای دوم}
%\secondsupervisorrank{استادیار}
% نام استاد(دان) مشاور را وارد کنید. چنانچه استاد مشاور ندارید، دستورات پایین را غیرفعال کنید.
%\firstadvisor{دکتر مشاور اول}
%\firstadvisorrank{استادیار}
%\secondadvisor{دکتر مشاور دوم}
% نام داوران داخلی و خارجی — نزدیک دفاع تکمیل می‌شود، فعلاً لازم نیست.
%\internaljudge{دکتر داور داخلی}
%\internaljudgerank{دانشیار}
%\externaljudge{دکتر داور خارجی}
%\externaljudgerank{دانشیار}
%\externaljudgeuniversity{دانشگاه داور خارجی}
% نام نماینده کمیته تحصیلات تکمیلی در دانشکده \ گروه — نزدیک دفاع تکمیل می‌شود.
%\graduatedeputy{دکتر نماینده}
%\graduatedeputyrank{دانشیار}
% نام دانشجو را وارد کنید
% TODO
\name{}
% نام خانوادگی دانشجو را وارد کنید
% TODO
\surname{}
% شماره دانشجویی دانشجو را وارد کنید
% TODO
\studentID{}
% تاریخ پایان‌نامه را وارد کنید
% TODO
\thesisdate{}
% به صورت پیش‌فرض برای پایان‌نامه‌های کارشناسی تا دکترا به ترتیب از عبارات «پروژه»، «پایان‌نامه» و «رساله» استفاده می‌شود؛ اگر  نمی‌پسندید هر عنوانی را که مایلید در دستور زیر قرار داده و آنرا از حالت توضیح خارج کنید.
%\projectLabel{پایان‌نامه}

% به صورت پیش‌فرض برای عناوین مقاطع تحصیلی کارشناسی تا دکترا به ترتیب از عبارت «کارشناسی»، «کارشناسی ارشد» و «دکتری» استفاده می‌شود؛ اگر نمی‌پسندید هر عنوانی را که مایلید در دستور زیر قرار داده و آنرا از حالت توضیح خارج کنید.
%\degree{}
%%%%%%%%%%%%%%%%%%%%%%%%%%%%%%%%%%%%%%%%%%%%%%%%%%%%
%% پایان‌نامه خود را تقدیم کنید! (اختیاری) %%
%% متن نمونه‌ی قالب اصلی (خانواده‌ی نگارنده‌ی قالب) غیرفعال شد — این صفحه اختیاری است؛
%% اگر می‌خواهید، متن خودتان را اینجا بنویسید و دستور \dedication را فعال کنید.
%\dedication
%{
%{\Large تقدیم به:}\\
%\begin{flushleft}{
%	\huge
%	...
%}
%\end{flushleft}
%}
%% متن قدردانی (اختیاری) %%
%% متن نمونه‌ی قالب اصلی غیرفعال شد — اگر می‌خواهید، متن خودتان (شامل نام استاد راهنما)
%% را اینجا بنویسید و دستور \acknowledgement را فعال کنید.
%\acknowledgement{
%...
%}
%%%%%%%%%%%%%%%%%%%%%%%%%%%%%%%%%%%%
%چکیده پایان‌نامه را وارد کنید
% TODO — چکیده‌ی فارسی (حداکثر ۱ صفحه) در فاز نگارش (پس از تکمیل فصل‌های ۲ تا ۵) نوشته
% می‌شود؛ نه پیش از آن. منبع محتوا: doc/project-definition.md + نتایج نهایی فصل ۴.
\fa-abstract{
}
% کلمات کلیدی پایان‌نامه را وارد کنید
% TODO تأیید یا اصلاح کنید
\keywords{پیش‌بینی تقاضا، رگرسیون کوانتایل، مسئله‌ی روزنامه‌فروش، سری‌های زمانی، هدررفت غذا}
% انتهای وارد کردن فیلد‌ها
%%%%%%%%%%%%%%%%%%%%%%%%%%%%%%%%%%%%%%%%%%%%%%%%%%%%%%
